\scanpage{39}

% B01 | printed letterhead, upper left
% letterhead
\begin{center}
UNIVERSITÉ DE PARIS\\
---$\cdot$---\\
ÉCOLE NORMALE SUPÉRIEURE\\
---$\cdot$---\\
45, RUE D'ULM\\
ODÉON 06$\cdot$45\\
---$\cdot$---
\end{center}

% B02 | cancelled line beside the address
\deleted{\unsure{En particulier si $E$ est}}

% B03 | Th 2, upper paragraph
\uuline{Th 2} : Si $G$ est un groupe fini d'ordre $n$, l'image de $p_*:H_*(E)\longrightarrow H_*(B)$ contient tous les éléments \deleted{d'ordre $n$} divisibles par $n$ dans $H_*(B)$.

% B04 | Dém, first two lines
\uuline{Dém} : Soit $\zeta$ \deleted{\unsure{un cycle}} $=\sum_{i\in I}a_i\sigma_i$ un cycle de $B$, où $\sigma_i$ sont des simplexes. Pour tout $\sigma_i$, soit $\tau_i$ simplexe de $E$ tel que $p\tau_i=\sigma_i$. Alors

% B05 | large sum and conclusion above fold
\[
\theta=\sum_{g\in G,\,i\in I}ga_i\tau_i\ \text{est un cycle de } E,\ \text{et }p(\theta)=\deleted{\unsure{\zeta|g}}n\zeta\ \text{\deleted{\unsure{est stable|est un}}}.
\]

% B06 | below right of sum
D'où le Th 2.

% B07 | Th 3, across and below fold
\uuline{Th 3} : Si \deleted{\unsure{$G$ est un|lorsque}} $E$ est \deleted{\unsure{acycle|acyclique}} contractile, $S={}$ ss-groupe engendré par les stabilisateurs fini d'ordre $n$, $\pi_i(B)$ est un groupe \deleted{\unsure{fini}} dont \inserted{tous les éléments sont} d'ordre qui divise une puissance de $n$ pour $i>1$

% B08 | Dém, last three lines
\uuline{Dém} : \deleted{\unsure{Comme|D'après}} D'après le Th. 1, $E/S$ est le revêtement universel de $B$, d'où $\pi_i(B)=\pi_i(E/S)$. D'après le Th. 2, \deleted{\unsure{par}} les $H_i(E/S)$

% archive sheet number: 38

